% !TeX root = ../papers/nc-p3.tex
\section{全正則窓の純関係スキームの有限性}
\label{sec:finite-pure-relation-locus}
\begin{theorem}\label{thm:finite-pure-relation-locus}
標準Koszul左右AS正則窓の各辺において、純テンソル関係のスキームは零次元以下である。
曲線射影に分岐があれば全ての順・逆中心の中央支持は3か4であり、
命題\ref{prop:finite-pure-pencil-opposite}の鉛筆と位相ごとの中心数上界が無条件に適用できる。
\end{theorem}
\begin{proof}[一文字を共有する純関係の排除]
曲面以上、有限辺、一次元の両側一点グラフでは、前定理の冒頭の理由で純関係はない。
曲面射影のない分岐の場合、一辺 $R_i$ が $u\otimes v,u\otimes v'$ を含み $v,v'$ は独立と仮定する。
平面 $H=\mathbb P(u^\perp)$ に制限して残る四式の四次行列式を $G$ とする。
$G$ は恒等零ではない。零なら全平面が射影像に入るからである。
$G=0$ の各既約成分の一般点で、重複度を保った平面上の局所環 $D=\mathcal O_{H,\eta}/(G)$ を取る。
順評価の三次小行列の一つは単元である。実際、階数2以下の点は有限個である。
ゆえに $D$ 上で行列の核は直和因子の直線で、一意な後点 $y$ が得られる。
その剰余体での値が動けば $y$ の順評価階数は3以上である。定点なら、
元の成分の全点がその前点なので逆中心であり、その順評価階数は3である。
したがって $D$ 上でも $y$ の順評価には階数3の直和因子がある。
隣接切断の完全対合と辺の消滅により元の点での逆評価はこの直交補に入る。
従って逆方向の全四次小行列式は $D$ 上で消える。この議論は $G$ の因子の全重複度を数えている。

平面上で逆四次小行列は全て $G$ の定数倍となり、少なくとも一つが非零である
（全て零なら逆曲面射影となる）。よって平面の一般点における
$R_{i-1}\to V_{i-1}$ の二次元核はPlücker座標が一定で、同じ二関係が平面全体で消える。
これらは $V_{i-1}\otimes ku$ に属し、$u\otimes v$ と連続する純関係を作って矛盾する。
逆向きも同じで、純関係の母線はどちらの向きにも存在しない。
\end{proof}

\begin{proof}[二次曲線と二項中心]
定理\ref{thm:pure-relation-curve-supports}により残る正次元の可能性は滑らかな二次曲線である。
その三関係は二空間 $A,B$ について $A\otimes B$ にあり、$A\otimes V$ に支持される
関係はこの三関係で尽きる。$L=\mathbb P(A^\perp)$ の一般点 $x$ で残り三関係を評価した
三列 $K_j(x)$ は独立である（可動中心の不在）。
$w(x)=\sum K_j(x)\otimes T_j$ と書けばその逆支持は $T_j$ の固定支持の和である。
後点が動けばその順階数は一般に3以上、後点が一定なら逆中心で順階数3なので、
直交により $x$ の逆階数は高々3である。一般に3であるから前点は一定 $p$ で、
$L$ がその後続直線である。命題\ref{prop:isolated-branch-blocks}による $p$ の三次元縮約像は
上記の三関係そのもので、中央支持は2である。

ここからは任意の中央支持2の順中心 $p$ を位置0に取って矛盾を示す。
定理\ref{thm:branch-binary-center-pairing}の対の逆中心を位置3の $z$、
その前点直線を位置2の $N$ とする。$p$ の後続直線を $L$ とする。
$p$ からの列の第二次数は1ではない。1なら次数公式から二直線列となり、
定理\ref{thm:two-line-full-middle-restriction}はその両端中央支持を3以上とするからである。

第二次数が2なら $L$ 上には唯一の順中心 $q$ があり、定理\ref{thm:rank-ten-conic-after-center}の
最初の標準二列の計算を用いる。この計算自体は反対階数を使わない。
独立パラメータでの制限は $[s,t]A(t)$、$A(t)$ は $q$ の鉛筆を、三行と中央支持の列で
書いた一次行列であり、第三点は $A(t)$ を消す。
$q$ は $L$ にあるので $w(q,z)=0$、その中央支持は3以下である。
中央支持2なら $A(t)$ は三行二列の $L_2$ で、左の核行は最小次数2である。
従って $[-,t]\otimes A(t)$ の列の全値を消す一次行はなく、これらは六空間全部を張る。
$R_0$ を $L$ に制限する写像が単射となるため、$A(t)$ の右核直線の全点が第二曲線の各一般点の後点となり、
可動中心になってしまう。中央支持3で常に特異なら、定数の左右核を除いた三次正方鉛筆は
$L_1\oplus L_1^{\mathsf t}$ に限る。一行二列のブロックが固定因子を共有する二純関係を与え、
前段に反する。正則三次正方鉛筆なら一般 $t$ の右核は四文字内で使わない定点一つであり、
二次曲線が終回になってしまう。従って第二次数は3である。
逆に $z$ からも第二次数は3となる。

$L$ 上の順欠損はこれで零である。$L$ の外の点の後点は二項対称二乗の三関係により
必ず $N$ に入る。よって位置1に順中心があればその直線全体が $N$ となり、次は $z$ で
動く位置が一つだけになってしまう。従って位置1には順中心がなく、逆に位置2には逆中心がない。
また二項中心の第二三次曲線の後の第三次数は4以上である。
定理\ref{thm:binary-centers-quadric-cones}の最初の錐の議論がここでも使える。
実際、三次から $N$ を中心とする射影の次数を $k>0$ とすれば第三曲線の頂点からの射影は
その二次写像で、$d_3\ge2k$。$d_3\le3$ なら $k=1$、三次曲線は $N$ と長さ2で会う。
第三曲線は二次または錐上の非平面三次で埋め込まれるが、その長さ2が $N$ の一定の後点 $z$ に写り矛盾する。
この議論は他の列に二直線列があるかによらない。逆向きにも同じ下界がある。

位置3に順中心があればその列の第二曲線は位置1にあり、一般前点が動くので $L$ と異なる。
二直線列なら位置2に逆中心を作り矛盾する。それ以外なら第三曲線は $N$ で、列はそこで動き終わり
$z$ に着く。これを $z$ から逆にたどる第三曲線は直線で、先の下界4に反する。
従って位置3にも順中心がない。残る順中心は偶位置だけで、次位相に順中心がないため
各列の第二次数は3である。これでは窓のどこにも二直線列がなく、
系\ref{cor:no-binary-centers-without-two-lines}が中央支持2自体を排除する。
逆中心も同様である。これで二次曲線もなくなり、有限性を得る。
\end{proof}
